\textbf{COLUNI 2018 -- Questão 33}\par\medskip O contexto atual de avanço tecnológico, configurando o espaço geográfico como meio técnico-científico-informacional, inaugurou a possibilidade de existência de empresas com fábricas globais. Nestas empresas a produção é descentralizada e espalhada por vários países: cada parte do processo produtivo é realizada num país. Ao final, o produto é destinado ao comércio internacional. A esse respeito, leia o fragmento a seguir:

\par

“Seus processos de produção são adaptados de modo a se aproveitar da disponibilidade de tecnologia e mão-de-obra, à medida que a empresa se estabelece. Seus métodos gerenciais são flexíveis e se amoldam às diversas culturas com que entra em contato. Sua cadeia de comando se baseia na tecnologia de informação e nas telecomunicações”. (Disponível em: https://midiaindependente.org/pt/red/2004/06/282690.shtml. Acesso em: 5 set. 2017)

\par

A configuração adotada na descentralização pelas empresas globais sinaliza o uso de vantagens locacionais específicas. NÃO é uma vantagem locacional dessa configuração descentralizada:\begin{enumerate}
\item Mão-de-obra barata em países subdesenvolvidos.
\item Incentivos e isenções fiscais de diversas naturezas.
\item Matéria prima obtida por meio de intensa regulação ambiental.
\item Transporte com técnicas que facilitam o fluxo de mercadoria.\end{enumerate}
